\scanpage{123}
On aura
\[
\begin{aligned}
\operatorname{Tr}u'&=\operatorname{Tr}u=b\\ \deleted{\unsure{\operatorname{Tr}}}\det u'&=\det u=c\\ \operatorname{Tr}u'v&=\operatorname{Tr}(b-u)v=b\underbrace{\operatorname{Tr}v}_{b'}-\underbrace{\operatorname{Tr}uv}_{\tau}=bb'-\tau=t
\end{aligned}
\]
et on gagne, en remplaçant $u$ par $u'$.

% slanted
\margin{Il manque le cas où $\delta\delta'=0$ cf plus bas pour ce cas-là} On a donc gagné dans le cas où on est sur un corps. Or on va maintenant voir que l'on a une réalisation \uline{canonique} de \deleted{$H,H'$} $(b,c,b',c',t)$ dans une alg. d'Azumaya, Pour ceci, on regarde l'alg.
% slanted
\margin{Attention au cercle vicieux, si on veut éviter la vérification fastidieuse de l'associativité de la loi de composition de $M(b,c,b',c',t)$}
\[
M(b,c,b',c',t)
\]
définie par ailleurs, à l'aide d'une base formelle
\[
\underline{1}\quad\underline{u},\ \underline{v},\ \underline{u}\,\underline{v}
\]
et une multiplication explicite. Tout revient à prouver que celle-ci est Azumaya. Il suffit de le voir sur les fibres géom -- on est ramené au cas d'un corps de base alg. clos. Mais là, l'argument
% archive sheet number: 122
